\chapter{Introduction}

\section{Background and Motivation}

Seed quality decides a large part of what a farmer harvests. A lot of seed that
looks fine to the eye can still carry a high fraction of broken, immature,
discoloured, or diseased grains, and that fraction sets the germination rate and
the final yield. Before a lot is planted, traded, or priced, someone has to judge
how good it is. In most of the supply chain that judgement is still made by a
human grader who pours out a sample, spreads it on a tray, and sorts the seeds by
hand into good and bad piles, counting as they go.

Manual inspection has two problems that are hard to fix by training people
better. It is slow: a careful grader works through a few hundred seeds before
fatigue sets in, which is a tiny sample of a multi-tonne lot. It is also
subjective. Two graders looking at the same tray will disagree on the borderline
seeds, and the same grader will drift over a long shift. The result is a quality
number that is neither repeatable nor auditable, which is a weak basis for a price
or a planting decision.

Computer vision fits this task well. The decision a grader makes, "is this
individual seed good or bad", is a per-object visual classification, and the wider
job of finding every seed in a tray is object detection. Both have matured to the
point where they run on commodity hardware and beat a tired human on consistency.
Surveys of deep learning in agriculture report strong results across crop and
produce inspection tasks once enough labelled data is
available~\cite{kamilaris2018deepag}, and the same convolutional models that
classify plant disease from leaf photographs transfer cleanly to grading the
seeds themselves~\cite{mohanty2016plant}. An automated grader does not get tired,
scores every seed in the photographed sample, and produces the same answer twice
for the same image.

This project, \textit{Seed Bank}, builds that automated grader as a working
platform rather than a notebook demo. It supports two crops today, coffee and
maize, chosen because they are economically important and visually distinct enough
to exercise the two-crop routing in the pipeline. A user photographs a batch of
seeds; the system locates each seed, scores its quality, and reports the aggregate
good-rate, seed count, confidence distribution, and a per-crop breakdown. One
backend serves both audiences: farmers in the field, through a web and a mobile
app, and the AI developers who run the models behind the product.

\section{Problem Definition}

Two linked problems sit at the centre of this project. The first is the shape of
the machine-learning task itself. The second is a data problem that the first one
exposes, and it is the harder of the two.

\subsection{Grading a tray photo is two ML problems in sequence}

A photo of a tray holds many seeds at once, so grading it cannot be a single
classification. It splits into two stages that run one after the other.

The first stage is \textbf{detection}: given the whole image, find where every
seed is. An object detector returns a bounding box, a confidence score, and a
class label for each seed it locates, turning one photo into $N$ located seeds. In
Seed Bank the detector's class identifies the crop, coffee or maize, because a
single tray can mix both and the next stage has to know which kind of seed it is
looking at.

The second stage is \textbf{classification}: for each detected seed, decide
whether it is good or bad. The system crops the bounding box out of the original
image and feeds that crop to a binary classifier that returns
\texttt{good}/\texttt{bad} plus a confidence. There is one classifier per crop
type, so a coffee crop goes to the coffee classifier and a maize crop to the maize
classifier. The data fans out along the way: one image becomes $N$ detections,
which become $N$ quality labels, and the aggregate report is computed from that
graph.

The two stages are genuinely different models with different training needs, and
keeping them separate is what lets a detector and a classifier be versioned and
swapped on their own. It also means the two stages need different training data,
which is where the second problem appears.

\subsection{The dataset gap}

The two stages do not have the same data appetite, and that asymmetry is the
problem this project had to solve before any of the platform work mattered.

The classifiers are easy to feed. A binary good/bad classifier trains on
\textbf{single-seed} crops, one seed centred in each image with a single label.
Public seed datasets are full of exactly this kind of image, because a single seed
on a plain background is cheap to photograph and trivial to label. Coffee and
maize both have usable single-seed corpora.

The detector is the problem. An object detector needs \textbf{multi-seed} images,
many seeds scattered in one frame, each one annotated with a bounding box. That
data barely exists in public. Almost every published seed dataset is single-seed,
which is useless for training a detector because there is nothing to locate, no
second object, no clutter, no overlap. Building a real multi-seed detection set by
hand means photographing trays of scattered seeds and drawing a box around every
grain in every image, hundreds of boxes per photo, thousands of photos. That
labelling effort is the bottleneck, and it is expensive and slow enough that it
stops most teams before they start.

So the situation is lopsided. The data needed to train the quality classifiers is
abundant, and the data needed to train the detector is scarce and costly to
produce. Closing that gap is what motivated the second half of this project,
\textit{MultiSeedGen}: a tool that takes the single-seed crops we already have and
synthesises annotated multi-seed images from them, so the detector can be trained
without hand-labelling thousands of trays.

\section{Project Objectives and Proposed Solution}

The objectives follow directly from the two problems above:

\begin{itemize}
  \item Build an end-to-end platform that takes a photo of a seed batch and
        returns a grading report, running the detect-then-classify pipeline for
        coffee and maize.
  \item Make every result traceable: any seed label must be linkable back to the
        exact model version and weights that produced it.
  \item Serve the same backend to three clients, a web app, a mobile app, and a
        programmatic API, so the product reaches both farmers and developers.
  \item Give the ML side real tooling: a model registry with a promotion
        lifecycle, a production-model resolver that selects the promoted
        production model for each segment, and offline evaluation of any model
        against labelled datasets.
  \item Close the detection-data gap by generating annotated multi-seed training
        images from single-seed crops.
\end{itemize}

The solution has two parts that match the two groups of objectives.

\textbf{Seed Bank} is the platform. Its backend is a FastAPI service that
orchestrates the two-stage pipeline: a request uploads images, a detector finds
the seeds, each seed is cropped and classified, and the results are persisted with
a hard foreign-key chain from \texttt{seed\_detection} to \texttt{inference} to
\texttt{model\_artifact}, which is what makes a result traceable. Around that
pipeline sits a model registry that moves weights through a
\texttt{registered} $\rightarrow$ \texttt{staging} $\rightarrow$ \texttt{production} $\rightarrow$ \texttt{archived}
lifecycle, a model resolver that selects the promoted production model for each
segment with a global fallback, and an experiment runner that scores models
offline against ground-truth datasets. Three clients consume one API: a React web
app, an Expo / React Native mobile app with a realtime camera, and the documented
HTTP API itself.

\textbf{MultiSeedGen} is the data generator. It cuts individual seeds out of
source photos, then composites many of them onto realistic backgrounds with
domain-matched lighting and camera degradation, and exports detection and
segmentation datasets in YOLO, YOLO-seg, and COCO formats with full provenance.
The cut-and-paste idea, pasting real object crops onto varied backgrounds to
synthesise detection training data, is a known and effective way to produce
labelled examples cheaply~\cite{dwibedi2017cutpaste}, and the deliberate
randomisation of lighting, scale, rotation, and noise follows the
domain-randomisation principle of forcing a model to generalise by training it on
a wide spread of synthetic variation~\cite{tobin2017domainrand}. Because every
seed is placed by the tool, its bounding box and mask are known exactly, so the
annotations come for free and the output is byte-reproducible for a fixed
configuration and random seed.

\section{Project Scope and Limitations}

\paragraph{In scope.} The platform grades two crops, coffee and maize, and is
built so a third crop is added by registering a new detector class and a matching
classifier rather than by rewriting the pipeline. It serves three clients (web,
mobile, API), resolves the promoted production model per segment with a global
fallback, runs offline evaluation of any registered model against labelled
datasets, and ships an observability stack covering structured logs, Prometheus
metrics, OpenTelemetry traces, and Sentry error reporting. MultiSeedGen covers the
full single-seed-to-multi-seed generation path with several segmentation backends
and reproducible output.

\paragraph{Out of scope and current limits.} The system is a working platform with
a known unfinished tail, recorded honestly so the reader is not misled:

\begin{itemize}
  \item \textbf{Upload.} Image upload to \texttt{POST /analyze} is multipart-only.
        Presigned direct-to-storage upload and batch cancellation were specified
        but never built.
  \item \textbf{Analytics ingestion.} The ClickHouse warehouse is populated by an
        application-level dual-write rather than true change-data-capture, even
        though Postgres is configured with \texttt{wal\_level=logical} so CDC can
        be added later.
  \item \textbf{Observability is opt-in.} OpenTelemetry tracing and Sentry are
        no-ops unless their endpoints are configured, so the default dev stack does
        not run them.
  \item \textbf{Some endpoints lack routes.} A few features have a service and a
        schema but no wired route yet (password change, experiment-results fetch,
        presigned upload).
  \item \textbf{The synthetic-to-real domain gap.} A detector trained only on
        MultiSeedGen output learns the statistics of synthetic scenes, which never
        match real trays perfectly. Narrowing that gap with augmentation and
        realistic degradation is the tool's central design concern, but it cannot
        be assumed closed.
\end{itemize}

\section{Development Methodology}

The platform was not written in one pass. It was delivered iteratively, in phases,
each one a usable slice of the system: scaffolding the FastAPI app, laying the
18-table async schema, wiring the infrastructure clients and health probes,
building auth, then the ML platform, then the inference path, experiments, the
data warehouse, and observability. That phasing is recorded in the project's
revamp-status document, and it meant each layer was exercised before the next was
stacked on it.

The backend follows clean-architecture principles~\cite{cleanarchitecture2017},
enforced as five rules rather than left as good intentions:

\begin{itemize}
  \item \textbf{Async end-to-end.} FastAPI, async SQLAlchemy, and async clients
        throughout, so nothing blocks the event loop.
  \item \textbf{Layered.} Every request flows
        \texttt{routers} $\rightarrow$ \texttt{services} $\rightarrow$ \texttt{repositories} $\rightarrow$ ORM.
        Routers stay thin, business logic lives in services, and only repositories
        touch the ORM.
  \item \textbf{Pydantic at the boundaries.} Every request and response is a
        validated schema; the ORM never leaks past the service layer.
  \item \textbf{One central \texttt{Settings}.} All configuration comes from a
        single settings object fed by environment variables or file secrets; code
        never reads the environment directly.
  \item \textbf{Model traceability.} The
        \texttt{seed\_detection} $\rightarrow$ \texttt{inference} $\rightarrow$ \texttt{model\_artifact}
        chain is enforced with foreign keys, not convention.
\end{itemize}

Work was managed with a git-based workflow. A pre-commit hook runs the formatter
and linter (ruff), the strict type checker (mypy), and a secret scanner (gitleaks)
before anything is committed, and the same checks are intended to run as CI gates.
The test suite is a pytest pyramid with unit, integration, and end-to-end tiers,
the integration tier using real Postgres and Redis through testcontainers. The
honest state of that suite, including the parts that are not yet hermetic, is
reported in the testing chapter.

\section{Tools and Technologies Used (Software and Hardware)}

The stack is deliberately conventional where it can be, so that each choice is a
well-understood tool with a known failure mode rather than a novelty. The backend
is built on FastAPI with async SQLAlchemy~\cite{fastapi,sqlalchemyguide}; model
experiments persist their metrics to PostgreSQL alongside a Markdown report
stored in MinIO, with no external tracking server; the web client is
React~\cite{reactindepth} and the mobile client is Expo / React
Native~\cite{reactnativeexpo}; and the whole system runs from a single Docker
Compose file~\cite{merkel2014docker}. Table~\ref{tab:ch1-tools} groups the full
set by area.

\begin{singlespace}
\begin{longtable}{p{0.22\textwidth} p{0.70\textwidth}}
  \caption{Tools and technologies used, grouped by area.}
  \label{tab:ch1-tools} \\
  \toprule
  \textbf{Area} & \textbf{Tools} \\
  \midrule
  \endfirsthead
  \toprule
  \textbf{Area} & \textbf{Tools} \\
  \midrule
  \endhead
  \bottomrule
  \endfoot
  Backend &
    Python 3.12, FastAPI, async SQLAlchemy~2 with asyncpg, Pydantic v2 and
    pydantic-settings, Alembic migrations. \\
  Async \& ML &
    Celery with a Redis broker, PyTorch and torchvision, Ultralytics YOLO,
    Roboflow inference SDK, Pillow and OpenCV (headless). \\
  Datastores &
    PostgreSQL 16 (OLTP), Redis 7 (cache, broker, rate-limit store), MinIO (object
    store), ClickHouse (analytics warehouse). \\
  Observability &
    structlog, Prometheus, Grafana, OpenTelemetry, Sentry. \\
  Web &
    React 18, TypeScript, Vite, Tailwind CSS, TanStack Query. \\
  Mobile &
    Expo SDK 52, React Native 0.76, expo-camera, React Navigation. \\
  Infrastructure &
    Docker and Docker Compose, a multi-stage Dockerfile, the \texttt{uv} dependency
    manager, and Make as the developer entrypoint. \\
\end{longtable}
\end{singlespace}

The split into many small runtime images is intentional. The API and the
lightweight worker carry no PyTorch at all, because torch and torchvision pull in
roughly 1.6\,GB and only the inference worker ever runs a model. That keeps the API
container small and the import cost of a routine background job low.

On the hardware side, development was done on a workstation with an NVIDIA RTX 3060
GPU. The GPU was used for two things: training the detector and classifier models,
and accelerating the rembg (U\textsuperscript{2}-Net) segmentation step in
MultiSeedGen when cutting seeds out of source photos. The Seed Bank platform itself
runs CPU-only in the development stack; the production Compose overlay swaps the
inference worker to a CUDA 12.4 GPU runtime with a device reservation, so the same
code path runs on a GPU in production without change.

\section{Project Timeline / Gantt Chart}

The work was scheduled across the 2025/2026 academic year and ran in the phased
order described in the methodology section above. The early part of the year went
to problem framing and to MultiSeedGen, because the platform's detector could not
be trained until synthetic detection data existed. The middle of the year built
the backend phase by phase, from the application scaffold and database schema
through auth, the ML platform, and the inference path. The later phases added
experiments, the analytics warehouse, and observability, in parallel with the web
and mobile clients and the bilingual localization work. The final stretch was
reserved for hardening, the production Compose overlay, and this report.
Figure~\ref{fig:gantt} shows the schedule.

\figslot{11-gantt.pdf}{Project timeline (Gantt chart).}{fig:gantt}

\section{Report Organization}

The rest of this report is organised as follows.

\textbf{Chapter 2, Related Work}, reviews the background this project builds on:
object detectors such as Faster R-CNN and YOLO and the feature-pyramid and
attention ideas behind them, image classifiers from the ResNet family, the
synthetic-data and augmentation literature that justifies MultiSeedGen, and prior
applications of computer vision in agriculture.

\textbf{Chapter 3, System Analysis}, sets out the requirements. It defines the user
roles, the functional and non-functional requirements, and the use cases the
platform has to support, and frames them against the two problems from this
chapter.

\textbf{Chapter 4, System Design}, describes the architecture: the layered backend,
the data model and its traceability chain, the two-stage inference pipeline with
its model registry and model resolver, the data warehouse, and the three clients.

\textbf{Chapter 5, Implementation}, covers how the design was actually built, the
model builders and inference backends, the Celery orchestration of the
detect-then-classify flow, MultiSeedGen's generation pipeline, and the deployment
setup.

\textbf{Chapter 6, Testing and Evaluation}, reports the test pyramid and quality
gates, the model evaluation results, and an honest account of the suite's current
gaps.

\textbf{Chapter 7, Conclusion and Future Work}, summarises what was delivered
against the objectives and lays out the unfinished tail and the next steps,
including true change-data-capture, the missing endpoints, and narrowing the
synthetic-to-real domain gap.
